% 第8.2节；首次阅读翻译，按用户要求不进行额外复核。
\section{例子与构造}\label{sec:8.2}
\setcounter{equation}{8}
\setcounter{figure}{9}
\subsection{凝聚层}\label{sec:8.2.1}

\SourceParagraph{ch08-8.2.1-001}{596}
这里先考虑一个最容易描述的 $A_\infty$ 范畴例子，并借此引出构造 $A_\infty$ 结构的一般方法。

\SourceParagraph{ch08-8.2.1-002}{596}
设 $X$ 为域 $k$ 上的非奇异代数簇，带有仿射开覆盖 $\mathcal U=\{U_i\}_{i\in I}$，其中 $I$ 为有序指标集。用 $D^b_\infty(X)$ 表示以局部自由层的有界复形为对象的范畴。为定义态射，先定义复形 $\mathcal{H}om^\bullet_{\mathcal O_X}(\mathcal E^\bullet,\mathcal F^\bullet)$，其中
\[
\mathcal{H}om^q_{\mathcal O_X}(\mathcal E^\bullet,\mathcal F^\bullet)
=\bigoplus_m\mathcal{H}om_{\mathcal O_X}(\mathcal E^m,\mathcal F^{m+q})
\]
由次数为 $q$ 的分次映射组成，而
\[
\delta\colon
\mathcal{H}om^q_{\mathcal O_X}(\mathcal E^\bullet,\mathcal F^\bullet)
\longrightarrow
\mathcal{H}om^{q+1}_{\mathcal O_X}(\mathcal E^\bullet,\mathcal F^\bullet)
\]
将 $(f_m\colon\mathcal E^m\to\mathcal F^{m+q})_m$ 映为
\[
(\delta(f))_m
=d_{\mathcal F}\circ f_m-(-1)^q f_{m+1}\circ d_{\mathcal E}
\colon\mathcal E^m\longrightarrow\mathcal F^{m+q+1}.
\]
接着，取 $\mathcal{H}om^\bullet_{\mathcal O_X}(\mathcal E^\bullet,\mathcal F^\bullet)$ 的 Čech 复形的全复形：
\[
\operatorname{Hom}^n_{D^b_\infty(X)}(\mathcal E^\bullet,\mathcal F^\bullet)
=\bigoplus_{p+q=n}
\check C^p\bigl(\mathcal U,
\mathcal{H}om^q_{\mathcal O_X}(\mathcal E^\bullet,\mathcal F^\bullet)\bigr).
\]

\SourceParagraph{ch08-8.2.1-003}{597}
这里
\[
\check C^p(\mathcal U,\mathcal G)
=\bigoplus_{i_0<\cdots<i_p}\Gamma(U_{i_0\cdots i_p},\mathcal G),
\qquad
U_{i_0\cdots i_p}=U_{i_0}\cap\cdots\cap U_{i_p}.
\]
在复形 $\operatorname{Hom}^\bullet_{D^b_\infty(X)}(\mathcal E^\bullet,\mathcal F^\bullet)$ 上，可以利用全微分
\[
\partial=d+(-1)^p\delta
\]
定义微分 $\partial$，其中 $d\colon\check C^p(\mathcal U,\mathcal G)\to\check C^{p+1}(\mathcal U,\mathcal G)$ 为通常的 Čech 微分，其定义为
\[
(d\alpha)_{i_0\cdots i_{p+1}}
=\sum_{j=0}^{p+1}(-1)^j
\alpha_{i_0\cdots\widehat{i_j}\cdots i_{p+1}}
\big|_{U_{i_0\cdots i_{p+1}}}.
\]

\SourceParagraph{ch08-8.2.1-004}{597}
这一复形与导出范畴中通常的态射概念之间，具有如下关系：
\[
H^i\bigl(\operatorname{Hom}^\bullet_{D^b_\infty(X)}
(\mathcal E^\bullet,\mathcal F^\bullet)\bigr)
=\operatorname{Ext}^i(\mathcal E^\bullet,\mathcal F^\bullet).
\]
见\BookSectionRef{sec:4.5}{4.5}。暂且不取上同调，而把 $D^b_\infty(X)$ 看成一个 $A_\infty$ 范畴。按照已经选定的符号约定，定义
\[
m_1(f)=(-1)^{\deg f}\partial f.
\]
复合的定义如下。设
\[
\begin{aligned}
f&=(f_{n,i_0\cdots i_p})\in
\check C^p\bigl(\mathcal U,
\mathcal{H}om^q_{\mathcal O_X}(\mathcal E^\bullet,\mathcal F^\bullet)\bigr),\\
g&=(g_{n',j_0\cdots j_{p'}})\in
\check C^{p'}\bigl(\mathcal U,
\mathcal{H}om^{q'}_{\mathcal O_X}(\mathcal F^\bullet,\mathcal G^\bullet)\bigr).
\end{aligned}
\]
则复合 $g\circ f$ 为
\[
\begin{aligned}
(g\circ f)_{n,i_0\cdots i_{p+p'}}
&=\left((-1)^{q'p}
 g_{n+q,i_0\cdots i_{p'}}\circ
 f_{n,i_{p'}\cdots i_{p+p'}}\right)
 \big|_{U_{i_0\cdots i_{p+p'}}},\\
(g\circ f)_{n,i_0\cdots i_{p+p'}}&\colon
\mathcal E^n\big|_{U_{i_0\cdots i_{p+p'}}}
\longrightarrow
\mathcal G^{n+q+q'}\big|_{U_{i_0\cdots i_{p+p'}}}.
\end{aligned}
\]
再定义
\[
m_2(g,f)=(-1)^{\deg f}g\circ f.
\]
容易检验\BookEqRef{eq:8.3}{8.3}成立，等价地，\BookEqRef{eq:8.4}{8.4}成立。此外，层态射的复合本来就满足结合律，因此可令 $m_n=0$，$n\geq3$，从而得到一个 $A_\infty$ 范畴。

\SourceParagraph{ch08-8.2.1-005}{597}
当然，这比 $A_\infty$ 范畴的条件更强：它实际上是一个 DG 范畴，即微分分次范畴，其高阶乘法映射全部为零。既然这里的态射尚不是我们想要的态射，$A_\infty$ 结构也并不有趣，那么它为何会是我们所需要的东西？

\SourceParagraph{ch08-8.2.1-006}{597}
答案来自以下代数结论，最初由 Kadeishvili\BookCite{ref279}提出。我们将给出它在范畴情形中的一个版本，这一版本归于 Kontsevich 和 Soibelman\BookCite{ref311}。

\SourceParagraph{ch08-8.2.1-007}{597}
\BookStatementNumber{7}
\begin{theorem}\label{thm:8.7}
给定 $A_\infty$ 代数 $A$，在 $H^*(A)$ 上存在一个满足 $m_1=0$ 的 $A_\infty$ 代数结构，以及 $A_\infty$ 代数的拟同构 $i\colon H^*(A)\to A$。
\end{theorem}

\SourceParagraph{ch08-8.2.1-008}{598}
下面对任意 $A_\infty$ 范畴 $\mathcal A$ 解释一个更一般的构造。

\SourceParagraph{ch08-8.2.1-009}{598}
对每一对 $X,Y\in\operatorname{Ob}(\mathcal A)$，需要选择一个投影算子，即次数为零的链映射
\[
\Pi\colon\operatorname{Hom}_{\mathcal A}(X,Y)
\longrightarrow\operatorname{Hom}_{\mathcal A}(X,Y),
\qquad\Pi^2=\Pi,
\]
以及次数为 $-1$ 的链同伦
\[
H\colon\operatorname{Hom}_{\mathcal A}(X,Y)
\longrightarrow\operatorname{Hom}_{\mathcal A}(X,Y)[-1],
\]
满足
\[
\Pi-1=m_1H+Hm_1.
\]
这特别说明，$\Pi$ 的像是一个与 $\operatorname{Hom}_{\mathcal A}(X,Y)$ 同伦等价的复形。我们特别感兴趣的情形，是 $\Pi$ 的像实际上同构于 $H^*(\operatorname{Hom}_{\mathcal A}(X,Y))$；不过，无论是否如此，构造都适用。由此得到新的 $A_\infty$ 范畴 $\Pi\mathcal A$，其对象与态射为
\[
\operatorname{Ob}(\Pi\mathcal A)=\operatorname{Ob}(\mathcal A),
\qquad
\operatorname{Hom}_{\Pi\mathcal A}(X,Y)
=\Pi\bigl(\operatorname{Hom}_{\mathcal A}(X,Y)\bigr).
\]
接着作如下定义。
\begin{enumerate}
\item $m_1^{\Pi\mathcal A}=m_1^{\mathcal A}$。由于 $\Pi$ 是链映射，此定义有意义。
\item $m_2^{\Pi\mathcal A}=\Pi\circ m_2^{\mathcal A}$。
\item $m_d^{\Pi\mathcal A}=\sum_Tm_{d,T}$，其中 $T$ 遍历具有 $d$ 个输入和一个输出的所有带状树，树的边上不附加长度。将 $m_n$ 配置到每个价数为 $n+1$ 的内部顶点，将 $H$ 配置到每条内部边，将 $\Pi$ 配置到输出顶点，并将 $m_{d,T}$ 定义为由该图自然给出的逐次复合。例如，对\BookRef{fig:8.10}{图10}所示的树 $T$，得到
\[
m_{d,T}(a_5,\ldots,a_1)
=\Pi\bigl(m_2(a_5,H(m_3(a_4,a_3,H(m_2(a_2,a_1)))))\bigr).
\]
\end{enumerate}

\SourceParagraph{ch08-8.2.1-010}{598}
此外，还可以定义一个 $A_\infty$ 函子 $i\colon\Pi\mathcal A\to\mathcal A$。在对象层面，$i$ 是恒等映射；$i_1$ 就是链复形的包含映射，而 $i_d$ 的定义与 $m_d$ 类似，写成
\[
i_d=\sum_Ti_{d,T},
\]
其中 $T$ 遍历相同的树，将输出处的 $\Pi$ 换成 $H$，就得到 $i_{d,T}$ 的定义。

\begin{figure}[H]
\centering
\BookImage{01_Numbered_Figures/Chapter_08/Ch08_Fig10_PDFp599_Bookp589.pdf}
\BookNumberOnlyCaption\label{fig:8.10}
\end{figure}

\SourceParagraph{ch08-8.2.1-011}{599}
下面说明这一构造为何成立。首先考虑 $m_2$ 的关系，记住 $m_1^{\Pi\mathcal A}=m_1$：
\[
\begin{aligned}
&m_2^{\Pi\mathcal A}(a_2,m_1(a_1))
+(-1)^{\deg a_1-1}m_2^{\Pi\mathcal A}(m_1(a_2),a_1)
+m_1\bigl(m_2^{\Pi\mathcal A}(a_2,a_1)\bigr)\\
&\quad=\Pi\bigl(m_2(a_2,m_1(a_1))
+(-1)^{\deg a_1-1}m_2(m_1(a_2),a_1)\bigr)
+m_1\bigl(\Pi(m_2(a_2,a_1))\bigr)\\
&\quad=0.
\end{aligned}
\]
这里使用了 $m_2$ 的关系；在最后一项中，还使用了 $\Pi$ 与 $m_1$ 可交换以及 $\Pi^2=\Pi$。

\SourceParagraph{ch08-8.2.1-012}{599}
检验 $d\geq3$ 时的 $m_d^{\Pi\mathcal A}$ 满足 $A_\infty$ 关系，并检验 $\{i_d\}$ 定义了一个 $A_\infty$ 函子，是一个很好的练习。下面按照\BookCite{ref311}概述论证，将逐一检验符号这一相当费力的工作留给读者。

\SourceParagraph{ch08-8.2.1-013}{599}
定义新映射
\[
\begin{aligned}
\widehat m_d^{\Pi\mathcal A}\colon{}&
\operatorname{Hom}_{\Pi\mathcal A}(X_{d-1},X_d)
\otimes\cdots\otimes
\operatorname{Hom}_{\Pi\mathcal A}(X_0,X_1)\\
&\longrightarrow\operatorname{Hom}_{\Pi\mathcal A}(X_0,X_d)[3-d],
\end{aligned}
\]
其定义为
\[
\widehat m_d^{\Pi\mathcal A}(a_d,\ldots,a_1)
=\sum_{(T,e)}\pm\widehat m_{(T,e)}(a_d,\ldots,a_1),
\]
其中求和遍历上述所有树 $T$，以及在 $T$ 中选定的一条边 $e$。$\widehat m_{(T,e)}$ 的定义与以前相同，但若 $e$ 是内部边，就将该边上的标记 $H$ 换为 $Hm_1+m_1H=\Pi-1$；若 $e$ 是外部边，就在这条边上加上一个 $m_1$，见\BookRef{fig:8.11}{图11}。每棵树必须选择一个依赖于 $e$ 的符号；我们将确定使下述论证成立的符号这一任务留给读者。

\begin{figure}[H]
\centering
\BookImage{01_Numbered_Figures/Chapter_08/Ch08_Fig11_PDFp600_Bookp590.pdf}
\BookNumberOnlyCaption\label{fig:8.11}
\end{figure}

\SourceParagraph{ch08-8.2.1-014}{600}
现在用两种不同的方法计算 $\widehat m_d^{\Pi\mathcal A}$。首先使用表达式 $\Pi-1$，于是得到对 $\widehat m_d^{\Pi\mathcal A}$ 的三类不同贡献：
\[
\widehat m_d^{\Pi\mathcal A}
=m_d^{\Pi\mathcal A,m_1}
+m_d^{\Pi\mathcal A,\Pi}
-m_d^{\Pi\mathcal A,1}.
\]
这里，$m_d^{\Pi\mathcal A,m_1}$ 由 $e$ 为外部边的各对 $(T,e)$ 所给出的贡献组成。其余两项通过对 $e$ 为内部边的各对 $(T,e)$ 求和定义：在这两种情形中，分别将边上的 $H$ 换为 $\Pi$ 和 $1$。符号因子与 $\widehat m_d^{\Pi\mathcal A}$ 定义中的符号因子相同。

\SourceParagraph{ch08-8.2.1-015}{600}
另一方面，可以在每条内部边上使用表达式 $m_1H+Hm_1$。这样，$e$ 为内部边的每一对 $(T,e)$ 都产生两个贡献。现在考虑任意内部顶点 $v$。对每条以 $v$ 为端点的边，$\widehat m_d^{\Pi\mathcal A}$ 中都恰有一项在该边紧邻 $v$ 的位置带有一个 $m_1$，如\BookRef{fig:8.12}{图12}所示。注意，外部边上不配置 $H$。因此，对这样的顶点，可以应用 $m_n^{\mathcal A}$ 的 $A_\infty$ 关系，将这个和写成\BookRef{fig:8.13}{图13}所示形式的各项之和；这里的符号实际上由 $m_n$ 的 $A_\infty$ 关系中的符号决定。由此可见，$\widehat m_d^{\Pi\mathcal A}$ 可以写为 $m_d^{\Pi\mathcal A,1}$ 中出现的各棵树的贡献之和。仔细处理符号，实际上得到
\[
\widehat m_d^{\Pi\mathcal A}=-m_d^{\Pi\mathcal A,1}.
\]
于是
\[
m_d^{\Pi\mathcal A,m_1}+m_d^{\Pi\mathcal A,\Pi}=0.
\]
这恰好就是 $m_d^{\Pi\mathcal A}$ 的 $A_\infty$ 关系。

\SourceParagraph{ch08-8.2.1-016}{600}
\begin{exercise*}
请看看这一论证对 $m_3$ 如何进行！
\end{exercise*}

\begin{figure}[H]
\centering
\BookImage{01_Numbered_Figures/Chapter_08/Ch08_Fig12_PDFp601_Bookp591.pdf}
\BookNumberOnlyCaption\label{fig:8.12}
\end{figure}
\begin{figure}[H]
\centering
\BookImage{01_Numbered_Figures/Chapter_08/Ch08_Fig13_PDFp601_Bookp591.pdf}
\BookNumberOnlyCaption\label{fig:8.13}
\end{figure}

\SourceParagraph{ch08-8.2.1-017}{601}
证明 $\{i_d\}_{d\geq1}$ 构成一个 $A_\infty$ 函子，论证是类似的。先检验次数为一和二的情形，再概述一般情形。对于 $d=1$，只需
\[
m_1^{\mathcal A}(i_1(a_1))=i_1(m_1^{\Pi\mathcal A}(a_1)),
\]
而 $i_1$ 是包含映射，所以这显然成立。对于 $d=2$，注意 $i_d$ 的定义给出
\[
i_2(a_2,a_1)=Hm_2^{\mathcal A}(a_2,a_1).
\]
考虑表达式
\[
\begin{aligned}
&Hm_2^{\mathcal A}(a_2,m_1(a_1))
+(-1)^{\deg a_1-1}Hm_2^{\mathcal A}(m_1(a_2),a_1)\\
&\hspace{5em}+Hm_1(m_2^{\mathcal A}(a_2,a_1)).
\end{aligned}
\]
利用 $m_2^{\mathcal A}$ 的关系，知道它实际上为零。不过，$Hm_1=-m_1H+\Pi-1$，因此，在前三项中使用 $i_2$ 的定义，在第四项中使用 $m_2^{\Pi\mathcal A}$ 的定义，就得到
\[
\begin{aligned}
0={}&i_2(a_2,m_1^{\Pi\mathcal A}(a_1))
+(-1)^{\deg a_1-1}i_2(m_1^{\Pi\mathcal A}(a_2),a_1)\\
&-m_1^{\mathcal A}(i_2(a_2,a_1))
+i_1m_2^{\Pi\mathcal A}(a_2,a_1)
-m_2^{\mathcal A}(i_1(a_2),i_1(a_1)),
\end{aligned}
\]
这就化为\BookEqRef{eq:8.8}{8.8}。

\SourceParagraph{ch08-8.2.1-018}{602}
一般情形的处理与前面基本相同。定义
\[
\widehat i_d=\sum_{(T,e)}\pm\widehat i_{(T,e)},
\]
求和仍遍历具有 $d$ 个输入的树，以及在树中选定的一条边；$\widehat i_{(T,e)}$ 的定义与 $\widehat m_{(T,e)}^{\Pi\mathcal A}$ 相同，只是将输出顶点所配置的算子 $\Pi$ 换成 $H$。符号的选择与 $\widehat m_d^{\Pi\mathcal A}$ 定义中的选择相同。在内部边上配置 $\Pi-1$，就可以写成
\[
\widehat i_d=i_d^{m_1}+i_d^\Pi-i_d^1,
\]
其中 $i_d^{m_1}$、$i_d^\Pi$ 和 $i_d^1$ 的定义与 $m_d$ 论证中完全相同，只是输出顶点上的 $\Pi$ 仍换成 $H$。若用 $m_1H+Hm_1$ 代替 $\Pi-1$ 来计算 $\widehat i_d$，就和前面一样得到
\[
\widehat i_d=-i_d^1,
\]
从而
\[
i_d^{m_1}+i_d^\Pi=0.
\]
现在可以看出，$i_d^\Pi$ 产生如下带有适当符号的和：
\[
\sum_{\substack{1<p<d\\0\leq q\leq d-p}}
\pm i_{d-p+1}\bigl(a_d,\ldots,a_{p+q+1},
 m_p^{\Pi\mathcal A}(a_{p+q},\ldots,a_{q+1}),
 a_q,\ldots,a_1\bigr).
\]
注意，这个和不包括 $p=1,d$ 的情形。$i_d^{m_1}$ 中选定的边为输入边的各项，给出 $p=1$ 的项：
\[
\sum_{\substack{p=1\\0\leq q\leq d-p}}
\pm i_{d-p+1}\bigl(a_d,\ldots,a_{q+2},
 m_1^{\Pi\mathcal A}(a_{q+1}),a_q,\ldots,a_1\bigr).
\]

\SourceParagraph{ch08-8.2.1-019}{602}
为得到余下的各项，考虑 $i_d^{m_1}$ 中选定的边为输出边的各项；它们给出形如 $Hm_1(\cdots)$ 的贡献。利用 $Hm_1=-m_1H+\Pi-1$，可见这些项中来自 $-m_1H$ 的贡献为 $-m_1^{\mathcal A}(i_d(a_d,\ldots,a_1))$，来自 $\Pi$ 的贡献为 $i_1m_d^{\Pi\mathcal A}(a_d,\ldots,a_1)$，而来自 $-1$ 的贡献稍微复杂一些：\BookRef{fig:8.14}{图14}所示的树贡献
\[
-m_r^{\mathcal A}\bigl(i_{s_r}(a_d,\ldots,a_{d-s_r+1}),
\ldots,i_{s_1}(a_{s_1},\ldots,a_1)\bigr).
\]
这涵盖了 $i_d$ 的关系中的全部项；再次仔细处理符号，就能得到 $i_d$ 的关系。

\begin{figure}[H]
\centering
\BookImage{01_Numbered_Figures/Chapter_08/Ch08_Fig14_PDFp603_Bookp593.pdf}
\BookNumberOnlyCaption\label{fig:8.14}
\end{figure}

\SourceParagraph{ch08-8.2.1-020}{603}
显然，$\{i_d\}_{d\geq1}$ 是 $A_\infty$ 范畴的拟同构，因为这只是对 $i_1$ 的一个条件，而链同伦 $H$ 的存在保证该条件成立。

\SourceParagraph{ch08-8.2.1-021}{603}
现在回到范畴 $D^b_\infty(X)$。上述构造的意义，在于能够给 $H^*(D^b_\infty(X))$ 配置一个 $A_\infty$ 结构。具体做法是：对每一对 $\mathcal E^\bullet,\mathcal F^\bullet$，选择复形 $\operatorname{Ext}^\bullet(\mathcal E^\bullet,\mathcal F^\bullet)$ 到 $\operatorname{Hom}^\bullet_{D^b_\infty(X)}(\mathcal E^\bullet,\mathcal F^\bullet)$ 中的包含映射 $i$，使其为拟同构；前一个复形在次数 $i$ 处为 $H^i(\operatorname{Hom}^\bullet_{D^b_\infty(X)}(\mathcal E^\bullet,\mathcal F^\bullet))$，边界映射为零。还要选择从 $\operatorname{Hom}^\bullet_{D^b_\infty(X)}(\mathcal E^\bullet,\mathcal F^\bullet)$ 投影到 $i(\operatorname{Ext}^\bullet(\mathcal E^\bullet,\mathcal F^\bullet))$ 的投影算子 $\Pi$ 和同伦 $H$。然后利用上述构造，在 $H^*(D^b_\infty(X))$ 上得到一个 $A_\infty$ 结构。这个分次范畴的零次上同调范畴就是通常的 $D^b(X)$，其各次态射则为导出范畴中的相应 $\operatorname{Ext}$ 群，所以这给出了 $D^b(X)$ 的一个 $A_\infty$ 结构。微分 $m_1$ 为零，但高阶乘法映射并不为零。此外，$D^b(X)$ 上的这一 $A_\infty$ 结构，给出一个与 $D^b_\infty(X)$ 拟同构的 $A_\infty$ 范畴。因此，尽管最初定义的高阶乘法映射为零，仍应将 $D^b_\infty(X)$ 上的 $A_\infty$ 结构看成一个极不平凡的结构。

\SourceParagraph{ch08-8.2.1-022}{603}
\BookStatementNumber{8}
\begin{example}\label{ex:8.8}
这里给出 $D^b(\mathbb P^1)$ 中非平凡 $m_3$ 的一个尽可能简单的例子，它出现于
\[
\begin{aligned}
m_3\colon{}&\operatorname{Ext}^1(\mathcal O_{\mathbb P^1}(2),\mathcal O_{\mathbb P^1})
\otimes\operatorname{Hom}(\mathcal O_{\mathbb P^1}(1),\mathcal O_{\mathbb P^1}(2))\\
&\otimes\operatorname{Hom}(\mathcal O_{\mathbb P^1},\mathcal O_{\mathbb P^1}(1))
\longrightarrow\operatorname{Hom}(\mathcal O_{\mathbb P^1},\mathcal O_{\mathbb P^1}).
\end{aligned}
\]

\SourceParagraph{ch08-8.2.1-023}{604}
为进行计算，需要下列层的 Čech 消解：
\[
\begin{aligned}
\mathcal{H}om(\mathcal O_{\mathbb P^1},\mathcal O_{\mathbb P^1}(1))
&=\mathcal{H}om(\mathcal O_{\mathbb P^1}(1),\mathcal O_{\mathbb P^1}(2))
=\mathcal O_{\mathbb P^1}(1),\\
\mathcal{H}om(\mathcal O_{\mathbb P^1},\mathcal O_{\mathbb P^1}(2))
&=\mathcal O_{\mathbb P^1}(2),\\
\mathcal{H}om(\mathcal O_{\mathbb P^1}(2),\mathcal O_{\mathbb P^1})
&=\mathcal O_{\mathbb P^1}(-2),\\
\mathcal{H}om(\mathcal O_{\mathbb P^1}(1),\mathcal O_{\mathbb P^1})
&=\mathcal O_{\mathbb P^1}(-1),\\
\mathcal{H}om(\mathcal O_{\mathbb P^1},\mathcal O_{\mathbb P^1})
&=\mathcal O_{\mathbb P^1}.
\end{aligned}
\]

\SourceParagraph{ch08-8.2.1-024}{604}
采用 $\mathbb P^1$ 的标准开覆盖 $U_0,U_1$，其中 $U_0=\mathbb C$，$U_1=(\mathbb C\setminus\{0\})\cup\{\infty\}$，$U_0$ 和 $U_1$ 上的坐标分别为 $z,z^{-1}$。下面给出这五个层的 Čech 复形，以及在需要时所选择的 $H$ 和 $\Pi$。在这种情形中，$H$ 仅在从次数一到次数零时非零。
\begin{table}[H]
\centering
\begingroup\small\setlength{\tabcolsep}{4pt}
\begin{tabular}{ll}
\toprule
$\mathcal O_{\mathbb P^1}(1)$
& $\mathbb C[z]\oplus z\mathbb C[z^{-1}]\longrightarrow\mathbb C[z,z^{-1}]$\\
$m_1$ & $(f_0,f_1)\longmapsto f_1-f_0$\\
\midrule
$\mathcal O_{\mathbb P^1}(2)$
& $\mathbb C[z]\oplus z^2\mathbb C[z^{-1}]\longrightarrow\mathbb C[z,z^{-1}]$\\
$m_1$ & $(f_0,f_1)\longmapsto f_1-f_0$\\
$\Pi_0$ & $(a_0+a_1z+a_2z^2+z^3f_0(z),z^2f_1(z^{-1}))$\\
& $\qquad\longmapsto(a_0+a_1z+a_2z^2,a_0+a_1z+a_2z^2)$\\
$\Pi_1$ & $0$\\
$H$ & $(z^3f_0(z)+z^2f_1(z^{-1}))\longmapsto(z^3f_0(z),-z^2f_1(z^{-1}))$\\
\midrule
$\mathcal O_{\mathbb P^1}(-2)$
& $\mathbb C[z]\oplus z^{-2}\mathbb C[z^{-1}]\longrightarrow\mathbb C[z,z^{-1}]$\\
$m_1$ & $(f_0,f_1)\longmapsto f_1-f_0$\\
\midrule
$\mathcal O_{\mathbb P^1}(-1)$
& $\mathbb C[z]\oplus z^{-1}\mathbb C[z^{-1}]\longrightarrow\mathbb C[z,z^{-1}]$\\
$m_1$ & $(f_0,f_1)\longmapsto f_1-f_0$\\
$\Pi$ & $0$\\
$H$ & $(f_0(z)+z^{-1}f_1(z^{-1}))\longmapsto(f_0(z),-z^{-1}f_1(z^{-1}))$\\
\midrule
$\mathcal O_{\mathbb P^1}$
& $\mathbb C[z]\oplus\mathbb C[z^{-1}]\longrightarrow\mathbb C[z,z^{-1}]$\\
$m_1$ & $(f_0,f_1)\longmapsto f_1-f_0$\\
$\Pi_0$ & $(a_0+zf_0(z),f_1(z^{-1}))\longmapsto(a_0,a_0)$\\
$\Pi_1$ & $0$\\
$H$ & $(zf_0(z)+f_1(z^{-1}))\longmapsto(zf_0(z),-f_1(z^{-1}))$\\
\bottomrule
\end{tabular}
\endgroup
\end{table}

\SourceParagraph{ch08-8.2.1-025}{604}
对 $m_3$ 有贡献的项只有三个，见\BookRef{fig:8.15}{图15}；不过，第一项无需考虑，因为在 Čech 复形层面没有 $m_3$。此外，
\[
\Pi m_2(a_3,Hm_2(a_2,a_1))=0,
\]
因为对 $a_2\in\operatorname{Hom}(\mathcal O_{\mathbb P^1}(1),\mathcal O_{\mathbb P^1}(2))$、$a_1\in\operatorname{Hom}(\mathcal O_{\mathbb P^1},\mathcal O_{\mathbb P^1}(1))$，$m_2(a_2,a_1)$ 的次数为零，而 $H$ 只在次数一上起作用。对于 $\Pi m_2(Hm_2(a_3,a_2),a_1)$，得到下表。
\begin{table}[H]
\centering
\begingroup\footnotesize\setlength{\tabcolsep}{3pt}
\begin{tabular}{cccccc}
\toprule
$a_3$ & $a_2$ & $a_1$ & $m_2(a_3,a_2)$
& $Hm_2(a_3,a_2)$ & $\Pi m_2(Hm_2(a_3,a_2),a_1)$\\
\midrule
$z^{-1}$ & $(1,1)$ & $(1,1)$ & $z^{-1}$ & $(0,-z^{-1})$ & $(0,0)$\\
$z^{-1}$ & $(z,z)$ & $(1,1)$ & $1$ & $(1,0)$ & $(1,1)$\\
$z^{-1}$ & $(1,1)$ & $(z,z)$ & $z^{-1}$ & $(0,-z^{-1})$ & $(0,0)$\\
$z^{-1}$ & $(z,z)$ & $(z,z)$ & $1$ & $(1,0)$ & $(0,0)$\\
\bottomrule
\end{tabular}
\endgroup
\end{table}

\SourceParagraph{ch08-8.2.1-026}{605}
最后一列描述了 $m_3$；特别地，它不为零。

\begin{figure}[H]
\centering
\BookImage{01_Numbered_Figures/Chapter_08/Ch08_Fig15_PDFp605_Bookp595.pdf}
\BookNumberOnlyCaption\label{fig:8.15}
\end{figure}
\end{example}
